%==============================================================================
% Trilingual abstract (English / French / Arabic).
%
% The Arabic page is typeset with ArabTeX (loaded in report-idris.tex), which
% keeps the whole report on pdflatex. Two rules apply inside `arabtext`:
%   * `\setcode{utf8}` must be active so UTF-8 Arabic input is read as Arabic;
%     it is switched back to `arabtex` afterwards so the rest of the document
%     is untouched.
%   * Latin runs (product names, acronyms, URLs) must be wrapped in `\LR{...}`,
%     otherwise their ASCII letters are read as ArabTeX transliteration.
% Parentheses are avoided in the Arabic text: ArabTeX does not mirror them for
% right-to-left, so they would render the wrong way round.
%==============================================================================

\fancyhead[R]{Abstract}
\chapter*{Abstract}
\addcontentsline{toc}{chapter}{Abstract}
\begin{spacing}{1.2}
%==============================================================================

In the European mid-sized enterprises that form
the economic backbone of DCH region (Germany, Austria and Switzerland), contracts are
still managed by hand, scattered across email inboxes, shared drives and
filing cabinets. Renewal and cancellation windows pass unnoticed, approvals
leave no auditable record, and the contract portfolio can be surveyed only
through a manual exercise that is stale the moment it is finished. This
happens at a time when GoBD, the European data-protection regulation and
ISO~27001 all expect business records to be retained immutably and traced to
an actor and a moment. The Contract Lifecycle Management (CLM) platforms that
could close this gap are scoped and priced for the large enterprise, or
require the customer to migrate every document into the vendor's own storage.

This report presents \textbf{CLMPilot}, an integration-first contract
lifecycle management platform for the DACH \textit{Mittelstand}, designed and
built during an end-of-studies internship at Gradient~Zero. Its defining
decision is to manage the lifecycle of a contract (its metadata, deadlines,
approvals and audit trail) while the document itself stays in the customer's
own storage, so adoption carries no migration cost. The platform is built on
\texttt{comby}, an in-house Go framework providing Command-Query
Responsibility Segregation and Event~Sourcing; that choice makes the
compliance-grade audit trail a property of the architecture itself rather
than a bolt-on feature.

\vspace{0.5cm}

\noindent\textbf{Keywords:} Contract Lifecycle Management, CQRS,
Event Sourcing, Domain-Driven Design, multi-tenancy, audit trail, Go,
Vue~3, DACH \textit{Mittelstand}.

\end{spacing}

%==============================================================================
% FRENCH
%==============================================================================
\fancyhead[R]{Résumé}
\chapter*{Résumé}
\begin{spacing}{1.2}

Dans les moyennes entreprises qui constituent l'épine dorsale économique de 
la region DACH (l'Allemagne, de
l'Autriche et de la Suisse), les contrats sont encore gérés manuellement et
dispersés entre boîtes de messagerie, disques partagés et armoires
d'archives. Les fenêtres de résiliation et de renouvellement passent
inaperçues, les validations ne laissent aucune trace auditable, et le
portefeuille contractuel ne peut être recensé que par un exercice manuel,
périmé dès qu'il est achevé. Or le GoBD, le règlement européen sur la
protection des données et la norme ISO~27001 exigent désormais que les
documents d'activité soient conservés de façon immuable et rattachés à un
acteur et à un instant. Les plateformes de gestion du cycle de vie des
contrats (CLM) susceptibles de combler ce manque sont dimensionnées et
tarifées pour les grands comptes, ou imposent au client de migrer l'ensemble
de ses documents vers le stockage de l'éditeur.

Ce rapport présente \textbf{CLMPilot}, une plateforme de gestion du cycle de
vie des contrats orientée intégration et destinée au \textit{Mittelstand} de
la région DACH, conçue et réalisée au cours d'un stage de fin d'études chez
Gradient~Zero. Le but fontamental de la plateforme est de gérer le cycle de vie du contrat 
(ses métadonnées, ses échéances, ses validations et sa piste d'audit) tandis
que le document lui-même reste dans le stockage du client : l'adoption
n'entraîne donc aucun coût de migration. La plateforme repose sur
\texttt{comby}, un framework Go interne mettant en {\oe}uvre la séparation
des responsabilités entre commandes et requêtes ainsi que l'\textit{Event
Sourcing} ; ce choix fait de la piste d'audit conforme aux exigences
réglementaires une propriété de l'architecture lui-même, et non une
fonctionnalité ajoutée.


\vspace{0.5cm}

\noindent\textbf{Mots-clés :} gestion du cycle de vie des contrats, CQRS,
Event Sourcing, conception pilotée par le domaine, multi-tenant, piste
d'audit, Go, Vue~3, \textit{Mittelstand} de la région DACH.

\end{spacing}

%==============================================================================
% ARABIC
%==============================================================================
% `\setcode{utf8}` is switched on here, at the top level and before the
% heading: inside a `\chapter*` argument that is still read under the default
% `arabtex` code, `\RL` swallows the title and re-emits it into the body.
\setcode{utf8}
\novocalize

\fancyhead[R]{Abstract (Arabic)}
\chapter*{\RL{ملخص}}
% ArabTeX's bidirectional line breaker fills a column about 4 lines short of
% the normal text height, which would push the last block onto a second page.
\enlargethispage{5\baselineskip}
\begin{spacing}{1.2}

% One `arabtext` per paragraph: a blank line inside the environment is read as
% an ordinary space, so paragraphs cannot be separated within a single block.
\begin{arabtext}
في الشركات الأوروبية المتوسطة الحجم التي تمثل العمود الفقري لاقتصاد منطقة \LR{DACH}
(ألمانيا والنمسا وسويسرا)، لا تزال العقود تدار يدويًا وتتوزع بين صناديق البريد
الإلكتروني والأقراص المشتركة وخزائن الملفات. ونتيجة لذلك، تفوت مواعيد التجديد
والإلغاء دون ملاحظتها، ولا تترك الموافقات أي أثر قابل للتدقيق، كما تعتري النظرة الشاملة
على محفظة العقود صعوبة بالغة إذ تتطلب جهدًا يدويًا يفقد صلاحيته بمجرد إنجازه.
يحدث هذا في وقت تشترط فيه قواعد \LR{GoBD} واللائحة الأوروبية لحماية البيانات ومعيار
\LR{ISO 27001} حفظ سجلات الأعمال بصورة غير قابلة للتعديل وتتبعها إلى المستخدم
واللحظة المحددة. أما منصات إدارة دورة حياة العقود (\LR{CLM}) القادرة على سد هذه
الفجوة، فإما أن تكون مصممة ومسعرة للمؤسسات الكبرى، أو تلزم العميل بترحيل جميع
وثائقه إلى بيئة التخزين الخاصة بالمزود.
\end{arabtext}

\begin{arabtext}
يقدم هذا التقرير منصة \LR{CLMPilot}، وهي منصة لإدارة دورة حياة العقود قائمة
على أولوية دمج الخصائص والوثائق وموجهة للشركات المتوسطة (\LR{Mittelstand}) في منطقة \LR{DACH}，
تم تصميمها وتطويرها خلال تدريب نهاية الدراسة في شركة \LR{Gradient Zero}. ويكمن
الهدف الأساسي للمنصة في إدارة دورة حياة العقد (بياناته الوصفية، وآجاله، وموافقاته،
وسجل تدقيقه) مع بقاء الوثيقة ذاتها في بيئة التخزين الخاصة بالعميل، مما يلغي تكاليف
الترحيل عند اعتمادها. تعتمد المنصة على \LR{comby}، وهو إطار عمل داخلي بلغة
\LR{Go} يطبق نمط فصل مسؤوليات الأوامر والاستعلامات (\LR{CQRS}) ونمط مصادر الأحداث
(\LR{Event Sourcing})؛ مما يجعل سجل التدقيق المطابق للمتطلبات التنظيمية خاصية أصيلة
في بنية النظام ذاته وليست مجرد ميزة إضافية.
\end{arabtext}

\vspace{0.3cm}

\begin{arabtext}
الكلمات المفتاحية: إدارة دورة حياة العقود، \LR{CQRS}، \LR{Event Sourcing}،
التصميم الموجه بالمجال، تعدد المستأجرين، سجل التدقيق، \LR{Go}،
\LR{Vue 3}، الشركات المتوسطة في منطقة \LR{DACH}.
\end{arabtext}

\setcode{arabtex}

\end{spacing}
